% !TEX root = ../../main.tex
\subsection{现有不足}

FoodShield 当前仍处于竞赛原型阶段，在以下方面存在不足：

第一，系统基于 SQLite 和单机后端完成核心流程验证，未在高并发通信、分布式存储和大规模订单数据处理场景下进行充分测试，系统吞吐量和并发能力有待评估。

第二，PID 映射仍由平台单点保存，存在平台单点信任风险。虽然系统通过审计日志和权限控制约束了管理员的溯源行为，但尚未引入盲签名、零知识证明或多方审批等去信任化机制。

第三，SM4-CBC 提供机密性保护，完整性由应用层 SM3 和 Merkle Tree 负责，两者独立运作。生产环境中可考虑采用认证加密模式（如 GCM），将加密与完整性验证合为一步，降低实现复杂度。

第四，系统当前以 HTTP 运行，通信链路未启用 TLS/WSS 加密。虽然竞赛演示环境下不影响安全机制验证，但在真实部署中需结合传输层安全措施。

第五，条件溯源的审批流程较为简化，尚未实现多级审批、多人联合授权和溯源理由模板等细粒度管控机制，管理员权限的约束仍依赖审计日志而非技术强制。
